\underline{定义} 2.7. \underline{若 $X$ 为 $\mathbb C$ 上的固有光滑代数簇, 则 $H^n(X)$ 的 Hodge 水平(Hodge level) 是满足 $p+q=n$ 且 $h^{p,q}(X)\ne0$ 的数 $q-p$ 的上确界.}

因此 Hodge 水平 $\ell$ 与 $n$ 模二同余, 且至多为
\[
\min\bigl(n,2\dim(X)-n\bigr).
\]
除 $H^n(X)=0$ 的情形外, 它非负; 在该情形下它为 $-\infty$.
\footnote{印本原文写作 $\min(n,\dim(X)-n)$. Hodge 指标的取值范围 $0\leq p,q\leq\dim(X)$ 要求因子 $2$; 英文版订正了原文这一明显的笔误.}

\underline{推论} 2.8. \underline{保留 2.1 的记号, 置 $k=\mathbb C$, 并假定 $a_d=\sup_i(a_i)$. 设 $\ell$ 为满足 $0\leq\ell<n$ 且与 $n$ 模二同余的整数. 则 $H^n(V_n(\underline a))$ 的 Hodge 水平至多为 $\ell$ 当且仅当}

%%%%%%%%%% source scan idx 66 | folio 59 %%%%%%%%%%

\begin{equation*}
\sum_{i\ne d}(a_i-1)
\leq (\ell+1)-(a_d-2)\frac{n-\ell}{2}.
\end{equation*}

由 (2.5.8), 若 Hodge 水平非负, 它就是零与下式中的最大值

\begin{equation*}
n-2\left\lfloor
\frac{n-\sum_{i\ne d}(a_i-1)+1}{a_d}
\right\rfloor.
\end{equation*}

不等式

\begin{equation*}
n-2\left\lfloor
\frac{n-\sum_{i\ne d}(a_i-1)+1}{a_d}
\right\rfloor\leq\ell
\end{equation*}

可改写为

\begin{equation*}
\frac{n-\ell}{2}
\leq
\left\lfloor
\frac{n-\sum_{i\ne d}(a_i-1)+1}{a_d}
\right\rfloor,
\end{equation*}

或等价地改写为

\begin{equation*}
\sum_{i\ne d}(a_i-1)
\leq 1+n-a_d\frac{n-\ell}{2}
=\ell+1-(a_d-2)\frac{n-\ell}{2}.
\end{equation*}

2.9. 利用公式 2.8, 对于小的 $\ell$, 不难列举 $\mathbb P^r$ 中维数为 $n>1$ 且其 $H^n$ 的水平为 $\ell$ 的完全交的多重次数. 然而, 本讲初版所给的表并不正确; 下表摘自 M. Rapoport 的文章 ``Complément à l'article de P. Deligne \emph{La conjecture de Weil pour les surfaces K3}'', \emph{Inventiones Math.} \underline{15}, 227--236 (1972) 中更完备的表. 当 $\ell\leq1$ 时, 维数为 $n$、多重次数为 $\underline a=(a_1,\ldots,a_d)$, 其中 $n>\ell$, $d\geq1$ 且 $a_i\geq2$, 并具有 Hodge 水平 $\ell$ 的完全交 $V_n(\underline a)$ 如下:

%%%%%%%%%% source scan idx 67 | folio 60 %%%%%%%%%%

\begin{equation*}
\begin{array}{ll}
\ell=-\infty & V_{2k+1}(2) \qquad (k\geq0),\\[2ex]
\ell=0 & V_{2k}(2)\ (k\geq1),\quad
          V_{2k}(2,2)\ (k\geq1),\quad V_2(3),\\[2ex]
\ell=1 & V_{2k+1}(2,2)\ (k\geq1),\quad
          V_{2k+1}(2,2,2)\ (k\geq1),\\
& V_3(3),\quad V_5(3),\quad V_3(3,2),\quad V_3(4).
\end{array}
\end{equation*}

2.10. 当 $\underline a$ 较小时, 由 2.4, $h_0^{p,q}(\underline a)$ 的生成函数为

\begin{equation*}
\begin{array}{rcl}
H(2)&=&\dfrac{1}{1-yz},\\[3ex]
H(3)&=&\dfrac{2+y+z}{1-3yz-y^2z-yz^2},\\[3ex]
H(2,2)&=&\dfrac{y+z}{(1-yz)^2}
          +\dfrac{2}{(1-yz)^2}+\dfrac{1}{1-yz},
\end{array}
\end{equation*}

所以

\begin{equation*}
\begin{array}{rcll}
h_0^{p,p}(2)&=&1&,\\[2ex]
h_0^{1,1}(3)&=&6&,\\[2ex]
h_0^{p,p}(2,2)&=&2p+3&.
\end{array}
\end{equation*}

%%%%%%%%%% source scan idx 68 | folio 61 %%%%%%%%%%

\bigskip
\begin{center}
\underline{参考文献}
\end{center}

\begin{enumerate}[label={[\arabic*]},leftmargin=3em]
\item Y. Akizuki 和 S. Nakano, ``Note on Kodaira--Spencer's proof of
Lefschetz theorems'', \emph{Proc. Jap. Acad.} 30, 266--272 (1954).

\item R. Hartshorne, \emph{Residues and Duality}, Lecture Notes in
Mathematics 20, Springer, 1966.

\item F. Hirzebruch, \emph{Topological Methods in Algebraic Geometry},
第 3 版, Springer-Verlag, 1966.

\item K. Kodaira, ``On a differential geometric method in the theory of
analytic stacks'', \emph{Proc. Nat. Acad. Sci. USA} 39, 1268--1273
(1953).
\end{enumerate}
